{\centering\textbf{Abstract}\par}\vspace{6pt}
\noindent We forecast next-day realized volatility for 95 NIFTY~100 stocks using 5-minute NSE data
(2016--2026) and ask whether deep learning beats the HAR benchmark, and whether the answer depends
on liquidity. On 29,774 out-of-sample stock-days (January 2025 -- April 2026), an LSTM fed with HAR
inputs (HAR-RV-LSTM) is the most accurate model: its QLIKE loss is 15.4\% below HAR-RV and 1.6\% below
log-HAR, both statistically significant. It beats HAR-RV on all 95 stocks and log-HAR on 70\% of them.
Its edge over log-HAR rises as liquidity falls, from 0.3\% (not significant) for the most liquid third
to 3.0\% for the least liquid third. An LSTM built on EGARCH features is the weakest model, and EGARCH
features add nothing to HAR-RV-LSTM. The gains come from calm and normal days. Correcting data errors,
the liquidity measure and look-ahead in the EGARCH features made the benchmark much stronger and the
neural models' advantage smaller than in our first version.

\medskip
\noindent\textbf{Keywords:} realized volatility, HAR, LSTM, EGARCH, liquidity, NSE, forecast evaluation\\
\textbf{JEL codes:} C45, C53, C58, G17
